\section{The full rank compression matroid lift}\label{sec:fulllift}
We prove Theorem~\ref{thm:main}. First suppose $Q>0$. The same construction includes $q=0$, with no original $B$ slots. Form a transversal matroid matched into the $q+2$ slots
\[
 B\sqcup\{e_1,e_2\}.
\]
Its elements and allowed slots are as follows:
\begin{itemize}
\item a private element $\bar b$ for each $b\in B$, allowed only in slot $b$;
\item an element for each $l\in L$, allowed in its original neighborhood in $B$;
\item two augmented core elements $\widehat A_i$, allowed in $S_i\cup\{e_i\}$;
\item two plane elements $P_1,P_2$, allowed only in $e_1,e_2$, respectively;
\item one common element $Z_j$ for each common right-exterior vertex, allowed in both $e_1,e_2$.
\end{itemize}
The rank is $q+2$, since all private $B$ elements together with $P_1,P_2$ form a basis. Every basis is counted once, regardless of how many injections into the slots it admits. We call $P_1,P_2,Z_j$ the plane elements, as they use only the two additional slots.

Give the private $B$ variables the values $s/v_b$, and the exterior-left variables $tu_l$. Keep core variables $X,Y$, and specialize the plane variables to
\begin{equation}\label{eq:planeweights}
 P_1:\frac{bs}{Q},\qquad P_2:\frac{as}{Q},\qquad
 Z_j:\frac{r_js}{Q}.
\end{equation}
Multiply this basis specialization by $Q\prod_{b\in B}v_b$. It is Lorentzian by the matroid and nonnegative-substitution theorems.

We now compute it exactly. Put
\[
 f_h=s^qf(t/s),\quad G_i=s^{q-1}g_i(t/s),\quad
 G_U=s^{q-1}g_U(t/s),\quad H_h=s^{q-2}h(t/s),
\]
where $G_i=G_U=0$ is defined directly for $q=0$, and $H_h=0$ is defined directly for $q<2$; the displayed homogenizing formulas are used only for nonnegative exponents. Classify a basis by the number $k$ of core elements and $\ell$ of plane elements. Other elements use only the $q$ original slots, so $k+\ell\ge2$ and $0\le k,\ell\le2$. There are six cases.
\begin{enumerate}
\item $(k,\ell)=(0,2)$. Any pair of distinct plane elements can fill the two extra slots, and the other $q$ elements fill the $B$ slots. The sum of plane-pair activities is
\[
 \frac{s^2}{Q^2}\left(ab+(a+b)c+\sum_{i<j}r_ir_j\right)=\frac{s^2}{Q}.
\]
After the exterior factor $Q$, the contribution is $s^2f_h$.
\item $(k,\ell)=(1,1)$. With core $\widehat A_1$, the plane element must be able to occupy $e_2$, hence is $P_2$ or a $Z_j$. Its total activity is $(a+c)s/Q$. The remaining $q$ elements fill the $B$ slots, giving $Xs(a+c)f_h$. With $\widehat A_2$, the contribution is $Ys(b+c)f_h$.
\item $(k,\ell)=(1,2)$. The plane pair fills both extra slots, forcing the core element into $B$. The remaining support uses row $S_i$ and $q-1$ ordinary elements, with profile $G_i$. The contribution is $Xs^2G_1+Ys^2G_2$.
\item $(k,\ell)=(2,0)$. The two core elements occupy their dedicated extra slots. The other $q$ elements fill $B$, contributing $QXYf_h$.
\item $(k,\ell)=(2,1)$. The element $P_1$ occupies $e_1$, forcing $\widehat A_2$ into $e_2$ and $\widehat A_1$ into $B$; hence its profile is $G_1$ and its activity coefficient is $b$. Similarly $P_2$ gives $G_2$ with coefficient $a$. A common $Z_j$ can occupy either extra slot, so feasibility holds exactly when either core row completes the remaining $B$-slot support. This is the union profile $G_U$, counted once even if both alternatives work. Thus the contribution is $XYs(bG_1+aG_2+cG_U)$.
\item $(k,\ell)=(2,2)$. The two plane elements fill the extra slots, forcing both core elements into $B$ together with $q-2$ other elements. The contribution is $XYs^2H_h$; this case is absent for $q<2$.
\end{enumerate}
Consequently the constructed Lorentzian polynomial is
\begin{align}\label{eq:fulllift}
 \mathcal L(s,t,X,Y)={}&s^2f_h+sX\big[(a+c)f_h+sG_1\big]
 +sY\big[(b+c)f_h+sG_2\big]\notag\\
 &+XY\big[Qf_h+s(bG_1+aG_2+cG_U)+s^2H_h\big].
\end{align}
Set $X=xt$ and $Y=yt$. Equation~\eqref{eq:union} identifies \eqref{eq:fulllift} with $s^{q+2}p_G(t/s)$, proving the bivariate assertion for $Q>0$. More generally, keep each exterior-left element variable as its own $z_l$ rather than $tu_l$, and set $X=z_{A_1}$, $Y=z_{A_2}$. The same six basis cases give exactly the homogeneous left marginal in Theorem~\ref{thm:main}. Every plane and private variable remains a nonnegative multiple of $s$, so this full marginal is Lorentzian. No analogous opposite-marginal assertion is inferred when $q>2$.

The six cases are disjoint classes of matroid bases. In particular the union case does not count injections or matching allocations. The construction also has a generic matrix representation obtained by giving each allowed slot incidence its own indeterminate. Its six block-determinant expansions provide a second interpretation of the same identity, including dependent or zero original core-row vectors.

If $Q=0$, add a singleton right vertex at each $A_i$ with activity $\varepsilon>0$. The perturbed graph still has a displayed $2+q$ cover, its $Q$ is positive, and its coefficients converge to those of the original graph as $\varepsilon\downarrow0$. Closedness gives the result. Other zero activities are handled by continuity at fixed order $q+2$; private reciprocals are used only before taking the positive-activity limit.

\begin{corollary}\label{cor:coverobstruction}
Any failure of weighted rank-normalized ultra-log-concavity must have every minimum vertex cover meet both shores in at least three vertices.
\end{corollary}
\begin{proof}
A pure or singleton-side minimum cover is covered by Lemma~\ref{lem:singleton}; a cover meeting one shore in exactly two vertices is covered by Theorem~\ref{thm:main}.
\end{proof}

\begin{remark}
The new rank-$(q+2)$ transversal matroid is used through its Lorentzian basis polynomial. The theorem asserts neither real-rootedness of $p_G$ nor stability of the original complemented-row polynomial. The support identity, and not the values of nonzero determinants, determines all coefficients.
\end{remark}
